\documentclass[a4paper,12pt]{article}
\usepackage[margin=1.5cm]{geometry}
\usepackage{amsmath, amssymb}
\usepackage{array}
\usepackage{alphalph}

\usepackage{amsmath, amssymb, amsfonts}
\usepackage{multicol}
\usepackage{multirow}

\usepackage{enumitem}
\usepackage{booktabs}
\usepackage{tabularx}
\usepackage{array}
\usepackage{framed}
\usepackage{xcolor}
\usepackage{tikz}
\usepackage{fancyhdr}

% Estilos para cajas de teoría y recordatorios
\definecolor{graybox}{gray}{0.92}
\newenvironment{boxsummary}{%
  \def\FrameCommand{\fcolorbox{black}{graybox}}%
  \MakeFramed{\FrameRestore}%
  \small\textbf{RECORDAR:}\par\smallskip
}{%
  \endMakeFramed
}

\begin{document}
\begin{enumerate}
    \item \textbf{Si $P(x) = x^2 - x + 2$ y $Q(x) = x^3 - x^2 + 1$, calcula:}
    \begin{multicols}{2}
    \begin{enumerate}[label=\textbf{\alph*)}]
        \item $P(1) + P(-1)$
        \item $P(0) + Q(-1)$
    \end{enumerate}
    \end{multicols}


    \item \textbf{¿Cuánto tiene que valer $a$ para que $P(a) = 0$ si $P(x) = 2x^2 - 3x + 1$?}
    \item \textbf{Realiza las siguientes divisiones de polinomios. Comprueba, en cada una de ellas, el resultado que obtienes.}
        \begin{enumerate}[label=\textbf{\alph*)}]
            \item $(2x^3 - 3x^2 - 5x - 5) : (x^2 - 2x - 1)$
            \item $(2x^3 - 3x^2 + 4x - 3) : (x^2 - 1)$
            \item $(x^4 + 1) : (x^2 + 1)$
            \item $(x^5 + 2x^3 - 1) : (x^2 - 3)$
        \end{enumerate}
    \item \textbf{El dividendo de una división de polinomios es $P(x) = x^5 - 2x^3 - x^2$, el cociente es $C(x) = x^2 - 2$ y el resto es $R(x) = -2$. ¿Cuál es el divisor?}
    \item \textbf{Determina el cociente y el resto, aplicando la regla de Ruffini.}
        \begin{enumerate}[label=\textbf{\alph*)}]
            \item $(x^3 - x^2 + x - 3) : (x - 1)$
            \item $(x^4 - x^3 - x + 9) : (x - 2)$
            \item $(x^4 + x^2 - 10) : (x - 5)$
            \item $(x^5 - 2x^3 + x - 7) : (x + 3)$
        \end{enumerate}
    \item \textbf{Calcula el valor de $m$ para que la división sea exacta.}
        \[
        (x^5 - 2x^3 - 8x^2 + mx + 3) : (x - 3)
        \]
    \item \textbf{Simplifica estas fracciones algebraicas.}

\begin{multicols}{2}
\begin{enumerate}[label=\textbf{\alph*)}]
    \item $\dfrac{2x - 2}{2x - 6}$
    
    \vspace{0.4cm}
    
    \item $\dfrac{x^2 - 1}{x^2 - 4x + 3}$
    
    \vspace{0.4cm}
    
    \item $\dfrac{x - 1}{x^2 - 1}$
    
    \columnbreak
    
    \item $\dfrac{2x^3 + 4x^2 + 2x}{6x^3 - 6x}$
    
    \vspace{0.4cm}
    
    \item $\dfrac{x^3 + 3x^2 - 4}{x^2 - 5x + 4}$
    
    \vspace{0.4cm}
    
    \item $\dfrac{x^2 - 1}{x^2 - 2x + 1}$
\end{enumerate}
\end{multicols}

\item  Pasar a álgebra las siguientes expresiones:
    \begin{enumerate}
        \item Un número “x” es 6 unidades mayor que la mitad de otro “y".
        \item La suma de las edades que tendrán una madre y una hija dentro de 10 años.
        \item El perímetro de un rectángulo es tres veces su altura.
        \item Dentro de 10 años la edad de mi padre será igual a la de mi madre actualmente.
        \item La diferencia de los inversos de dos números disminuida en 10 unidades.
        \item El cociente de dos números es 8.
        \item El cociente de dos números es 6 y el resto 4.
        \item La suma de los cuadrados de dos números consecutivos.
        \item El producto de dos números es 143 veces su cociente.
    \end{enumerate}



\end{enumerate}

Soluciones

\begin{enumerate}
    \item \textbf{Si $P(x) = x^2 - x + 2$ y $Q(x) = x^3 - x^2 + 1$, calcula:}
    \begin{enumerate}[label=\textbf{\alph*)}]
        \item $P(1) + P(-1) = (1 - 1 + 2) + (1 + 1 + 2) = 2 + 4 = 6$
        \item $P(0) + Q(-1) = 2 + (-1 - 1 + 1) = 2 + (-1) = 1$
    \end{enumerate}

    \vspace{0.5cm}

    \item \textbf{¿Cuánto tiene que valer $a$ para que $P(a) = 0$ si $P(x) = 2x^2 - 3x + 1$?}

    Son las soluciones de la ecuación $2x^2 - 3x + 1 = 0 \to x = 1$ y $x = \dfrac{1}{2}$
    \item \textbf{Realiza las siguientes divisiones de polinomios. Comprueba, en cada una de ellas, el resultado que obtienes.}

        \begin{enumerate}[label=\textbf{\alph*)}]
            \item $(2x^3 - 3x^2 - 5x - 5) : (x^2 - 2x - 1) \longrightarrow \begin{cases} \text{Cociente} = 2x + 1 \\ \text{Resto} = -x - 4 \end{cases}$

            \vspace{0.3cm}

            \item $(2x^3 - 3x^2 + 4x - 3) : (x^2 - 1) \longrightarrow \begin{cases} \text{Cociente} = 2x - 3 \\ \text{Resto} = 6x - 6 \end{cases}$

            \vspace{0.3cm}

            \item $(x^4 + 1) : (x^2 + 1) \longrightarrow \begin{cases} \text{Cociente} = x^2 - 1 \\ \text{Resto} = 2 \end{cases}$

            \vspace{0.3cm}

            \item $(x^5 + 2x^3 - 1) : (x^2 - 3) \longrightarrow \begin{cases} \text{Cociente} = x^3 + 5x \\ \text{Resto} = 15x - 1 \end{cases}$
        \end{enumerate}
    \item \textbf{El dividendo de una división de polinomios es $P(x) = x^5 - 2x^3 - x^2$, el cociente es $C(x) = x^2 - 2$ y el resto es $R(x) = -2$. ¿Cuál es el divisor?}

        \vspace{0.3cm}

        $P(x) = Q(x) \cdot C(x) + R(x)$

        $x^5 - 2x^3 - x^2 = Q(x) \cdot (x^2 - 2) - 2 \longrightarrow$
        \begin{align*}
        &\longrightarrow x^5 - 2x^3 - x^2 + 2 = Q(x) \cdot (x^2 - 2) \longrightarrow \\
        &\longrightarrow Q(x) = (x^5 - 2x^3 - x^2 + 2) : (x^2 - 2) = x^3 - 1
        \end{align*}
\item \textbf{Determina el cociente y el resto, aplicando la regla de Ruffini.}

\begin{enumerate}[label=\textbf{\alph*)}]
    \item $(x^3 - x^2 + x - 3) : (x - 1)$
    
    \begin{center}
    \begin{tabular}{r|rrrr}
      & 1 & $-1$ & 1 & $-3$ \\
    1 &   & 1 & 0 & 1 \\ \hline
      & 1 & 0 & 1 & \multicolumn{1}{|r}{$-2$}
    \end{tabular}
    \quad $\longrightarrow \quad C(x) = x^2 + 1; \quad R(x) = -2$
    \end{center}

    \vspace{0.4cm}

    \item $(x^4 - x^3 - x + 9) : (x - 2)$
    
    \begin{center}
    \begin{tabular}{r|rrrrr}
      & 1 & $-1$ & 0 & $-1$ & 9 \\
    2 &   & 2 & 2 & 4 & 6 \\ \hline
      & 1 & 1 & 2 & 3 & \multicolumn{1}{|r}{15}
    \end{tabular}
    \quad $\longrightarrow \quad C(x) = x^3 + x^2 + 2x + 3; \quad R(x) = 15$
    \end{center}

    \vspace{0.4cm}

    \item $(x^4 + x^2 - 10) : (x - 5)$
    
    \begin{center}
    \begin{tabular}{r|rrrrr}
      & 1 & 0 & 1 & 0 & $-10$ \\
    5 &   & 5 & 25 & 130 & 650 \\ \hline
      & 1 & 5 & 26 & 130 & \multicolumn{1}{|r}{640}
    \end{tabular}
    
    \smallskip
    $C(x) = x^3 + 5x^2 + 26x + 130; \quad R(x) = 640$
    \end{center}

    \vspace{0.4cm}

    \item $(x^5 - 2x^3 + x - 7) : (x + 3)$
    
    \begin{center}
    \begin{tabular}{r|rrrrrr}
      & 1 & 0 & $-2$ & 0 & 1 & $-7$ \\
    $-3$ &   & $-3$ & 9 & $-21$ & 63 & $-192$ \\ \hline
      & 1 & $-3$ & 7 & $-21$ & 64 & \multicolumn{1}{|r}{$-199$}
    \end{tabular}
    
    \smallskip
    $C(x) = x^4 - 3x^3 + 7x^2 - 21x + 64; \quad R(x) = -199$
    \end{center}

    \vspace{0.4cm}

    \item $(x^7 + x^4 - 7x^2) : (x + 4)$
    
    \begin{center}
    \small
    \begin{tabular}{r|rrrrrrrr}
      & 1 & 0 & 0 & 1 & 0 & $-7$ & 0 & 0 \\
    $-4$ &   & $-4$ & 16 & $-64$ & 252 & $-1.008$ & 4.060 & $-16.240$ \\ \hline
      & 1 & $-4$ & 16 & $-63$ & 252 & $-1.015$ & 4.060 & \multicolumn{1}{|r}{$-16.240$}
    \end{tabular}
    
    \smallskip
    $C(x) = x^6 - 4x^5 + 16x^4 - 63x^3 + 252x^2 - 1.015x + 4.060; \quad R(x) = -16.240$
    \end{center}
\end{enumerate}
\item \textbf{Calcula el valor de $m$ para que la división sea exacta.}
\[
(x^5 - 2x^3 - 8x^2 + mx + 3) : (x - 3)
\]

\begin{center}
\begin{tabular}{r|rrrrrr}
  & 1 & 0 & $-2$ & $-8$ & $m$ & 3 \\
3 &   & 3 & 9 & 21 & 39 & $117 + 3m$ \\ \hline
  & 1 & 3 & 7 & 13 & $39 + m$ & \multicolumn{1}{|r}{$120 + 3m$}
\end{tabular}
\end{center}

\vspace{0.3cm}

\[
120 + 3m = 0 \longrightarrow 3m = -120 \longrightarrow m = -40
\]

\item \textbf{Simplifica estas fracciones algebraicas.}

\begin{enumerate}[label=\textbf{\alph*)}]
    \item $\dfrac{2x - 2}{2x - 6} = \dfrac{2 \cdot (x - 1)}{2 \cdot (x - 3)} = \dfrac{x - 1}{x - 3}$
    
    \vspace{0.3cm}
    
    \item $\dfrac{x^2 - 1}{x^2 - 4x + 3} = \dfrac{(x + 1) \cdot (x - 1)}{(x - 1) \cdot (x - 3)} = \dfrac{x + 1}{x - 3}$
    
    \vspace{0.3cm}
    
    \item $\dfrac{x - 1}{x^2 - 1} = \dfrac{x - 1}{(x + 1) \cdot (x - 1)} = \dfrac{1}{x + 1}$
    
    \vspace{0.3cm}
    
    \item $\dfrac{2x^3 + 4x^2 + 2x}{6x^3 - 6x} = \dfrac{2x \cdot (x + 1)^2}{6x \cdot (x + 1) \cdot (x - 1)} = \dfrac{x + 1}{3 \cdot (x - 1)}$
    
    \vspace{0.3cm}
    
    \item $\dfrac{x^3 + 3x^2 - 4}{x^2 - 5x + 4} = \dfrac{(x - 1) \cdot (x + 2)^2}{(x - 4) \cdot (x - 1)} = \dfrac{(x + 2)^2}{x - 4}$
    
    \vspace{0.3cm}
    
    \item $\dfrac{x^2 - 1}{x^2 - 2x + 1} = \dfrac{(x - 1) \cdot (x + 1)}{(x - 1)^2} = \dfrac{x + 1}{x - 1}$
\end{enumerate}

\item \item \textbf{Pasar a álgebra las siguientes expresiones:}
\begin{enumerate}[label=\textbf{\alph*)}]
    \item Un número “x” es 6 unidades mayor que la mitad de otro “y”.
    \[
    x = \frac{y}{2} + 6
    \]

    \item La suma de las edades que tendrán una madre y una hija dentro de 10 años (siendo $m$ la edad de la madre e $h$ la de la hija).
    \[
    (m + 10) + (h + 10) = m + h + 20
    \]

    \item El perímetro de un rectángulo es tres veces su altura (siendo $b$ la base y $h$ la altura).
    \[
    2b + 2h = 3h \quad \text{o bien} \quad 2b = h
    \]

    \item Dentro de 10 años la edad de mi padre será igual a la de mi madre actualmente (siendo $p$ la edad del padre y $m$ la de la madre).
    \[
    p + 10 = m
    \]

    \item La diferencia de los inversos de dos números ($x$ e $y$) disminuida en 10 unidades.
    \[
    \frac{1}{x} - \frac{1}{y} - 10
    \]

    \item El cociente de dos números ($x$ e $y$) es 8.
    \[
    \frac{x}{y} = 8
    \]

    \item El cociente de dos números es 6 y el resto 4 (siendo $D$ el dividendo y $d$ el divisor, $D = d \cdot c + r$).
    \[
    D = 6d + 4 \quad \text{o bien} \quad \frac{D}{d} = 6 + \frac{4}{d}
    \]

    \item La suma de los cuadrados de dos números consecutivos ($x$ y $x+1$).
    \[
    x^2 + (x + 1)^2
    \]

    \item El producto de dos números es 143 veces su cociente.
    \[
    x \cdot y = 143 \cdot \frac{x}{y}
    \]
\end{enumerate}



\end{enumerate}

\end{document}

